% Math / Physics
% Look: numbered statements with an inline sans-serif label,
%   "Definition 3 (Homotopy). Let f, g : X -> Y be ..."
% in tinted boxes with a coloured bar on the left.
\newif\ifnotesframed
% For plain framed boxes instead, uncomment the next line.
% \notesframedtrue
\usepackage{amsthm}
\usepackage{thmtools}
\usepackage{siunitx}
\usepackage{bm}

\newcommand{\dd}{\mathop{}\!\mathrm{d}}
\renewcommand{\vec}[1]{\bm{#1}}
\DeclareMathOperator{\tr}{tr}
\DeclareMathOperator{\im}{im}
\DeclareMathOperator{\id}{id}

\definecolor{notesdef}{RGB}{38,128,82}
\definecolor{notesthm}{RGB}{150,78,36}
\definecolor{notesex}{RGB}{52,84,150}
\definecolor{notesrem}{RGB}{100,100,100}
\colorlet{notesproof}{notesthm!45}

% #1 style name, #2 colour of the label, #3 space around it (boxed
% statements get their spacing from the box instead)
\newcommand{\notesstyle}[3]{%
  \declaretheoremstyle[headfont=\sffamily\bfseries\color{#2}, notefont=\sffamily\mdseries\color{#2},
    notebraces={(}{)}, bodyfont=\normalfont, headpunct={.}, postheadspace=0.5em,
    spaceabove=#3, spacebelow=#3]{#1}}
\notesstyle{notesdefstyle}{notesdef}{0pt}
\notesstyle{notesthmstyle}{notesthm}{0pt}
\notesstyle{notesexstyle}{notesex}{8pt}
\notesstyle{notesremstyle}{notesrem}{8pt}

\declaretheorem[style=notesdefstyle, name=Definition]{definition}
\declaretheorem[style=notesthmstyle, name=Theorem]{theorem}
\declaretheorem[style=notesthmstyle, name=Lemma, sibling=theorem]{lemma}
\declaretheorem[style=notesthmstyle, name=Proposition, sibling=theorem]{proposition}
\declaretheorem[style=notesthmstyle, name=Corollary, sibling=theorem]{corollary}
\declaretheorem[style=notesexstyle, name=Example, numbered=no]{example}
\declaretheorem[style=notesremstyle, name=Remark, numbered=no]{remark}
\declaretheorem[style=notesremstyle, name=Note, numbered=no]{note}

\renewenvironment{proof}[1][\proofname]%
  {\par\noindent{\sffamily\bfseries\color{notesthm}#1.}\ }%
  {\hfill$\square$\par}

% Put each statement in a box. #1 environment, #2 colour
\newcommand{\notesbox}[2]{%
  \ifnotesframed
    \tcolorboxenvironment{#1}{breakable, colback=white, colframe=black!60,
      boxrule=0.5pt, sharp corners, left=6pt, right=6pt, top=4pt, bottom=4pt,
      before skip=10pt, after skip=10pt}%
  \else
    \tcolorboxenvironment{#1}{enhanced, breakable, frame hidden, sharp corners,
      colback=#2!7, borderline west={2pt}{0pt}{#2},
      left=8pt, right=6pt, top=5pt, bottom=5pt, before skip=10pt, after skip=10pt}%
  \fi}
\notesbox{definition}{notesdef}
\notesbox{theorem}{notesthm}
\notesbox{lemma}{notesthm}
\notesbox{proposition}{notesthm}
\notesbox{corollary}{notesthm}
\notesbox{proof}{notesproof}
\notesbox{note}{notesrem}
